\section*{अध्याय \ref{ch:g4:world-physical-map} --- विश्व का भौतिक मानचित्र}

\begin{solution}{exo:g4:world-physical-map:1}
किसी जगह की \omterm{def:g4:world-physical-map:altitude}{ऊँचाई} \omterm{def:g4:world-physical-map:altitude}{समुद्र तल} से उसकी \omterm{def:g4:world-physical-map:altitude}{ऊँचाई} है, मीटर में। वह \omterm{def:g4:world-physical-map:altitude}{समुद्र तल} से नापी
जाती है, जो \omterm{def:g2:water-on-earth:ocean}{समुद्र} की सतह का औसत तल है।
\end{solution}

\begin{solution}{exo:g4:world-physical-map:2}
(क) गहरा हरा: 0 और 200 मीटर के बीच। (ख) हल्का पीला: 500 और 1\,000 मीटर के
बीच। (ग) लगभग सफ़ेद: 4\,000 मीटर से ऊपर।
\end{solution}

\begin{solution}{exo:g4:world-physical-map:3}
(क) दक्षिणी अमेरिका। (ख) एशिया। (ग) अफ़्रीका। (घ) उत्तरी अमेरिका।
\end{solution}

\begin{solution}{exo:g4:world-physical-map:4}
सबसे ऊँचा \omterm{def:g1:landscapes-around-us:landform}{पर्वत} माउंट एवरेस्ट है (8\,849 मीटर)। सबसे बड़ी \omterm{def:g1:landscapes-around-us:water}{झील} कैस्पियन सागर है;
सबसे गहरी \omterm{def:g1:landscapes-around-us:water}{झील} बैकाल \omterm{def:g1:landscapes-around-us:water}{झील} है।
\end{solution}

\begin{solution}{exo:g4:world-physical-map:5}
नीला जितना गहरा, \omterm{def:g2:water-on-earth:ocean}{समुद्र} उतना गहरा: तटों के पास \omterm{def:g2:water-on-earth:ocean}{समुद्र} उथला है (200 मीटर से कम)
और \omterm{def:g2:water-on-earth:ocean}{महासागरों} के बीच में बहुत ज़्यादा गहरा (कई जगहों पर 4\,000 मीटर से
ज़्यादा)।
\end{solution}

\begin{solution}{exo:g4:world-physical-map:6}
\omterm{def:g4:world-physical-map:physical-map}{भौतिक मानचित्र} सतह की \omterm{def:g4:world-physical-map:altitude}{ऊँचाई} दिखाता है। अंटार्कटिका और ग्रीनलैंड में सतह बर्फ़ की
एक मोटी परत की ऊपरी सतह है, जो \omterm{def:g4:world-physical-map:altitude}{समुद्र तल} से बहुत ऊपर है, इसलिए उन्हें ऊँची ज़मीन
के रंग मिलते हैं।
\end{solution}

\begin{solution}{exo:g4:world-physical-map:7}
चैलेंजर गर्त का तल (\omterm{def:g4:world-physical-map:altitude}{समुद्र तल} से लगभग 11\,000 मीटर नीचे); \omterm{def:g2:water-on-earth:ocean}{महासागर} का औसत तल
(लगभग 3\,700 मीटर नीचे); एक \omterm{def:g2:water-on-earth:ocean}{समुद्र} \omterm{def:g1:landscapes-around-us:water}{तट} (लगभग 0 मीटर); 1\,500 मीटर पर एक गाँव;
माउंट एवरेस्ट की चोटी
(8\,849 मीटर)।
\end{solution}

\begin{solution}{exo:g4:world-physical-map:8}
एक \omterm{def:g1:landscapes-around-us:water}{नदी} बहुत-सी धाराओं से बनती है जो आपस में मिलती हैं, और लोग इस पर सहमत नहीं
कि कौन-सी धारा उसकी असली शुरुआत है। किसी दूर के स्रोत से नापने पर \omterm{def:g1:landscapes-around-us:water}{नदी} ज़्यादा
लंबी निकलती है। नील और अमेज़न की लंबाई इतनी पास-पास है (हर एक लगभग 6\,700 से
6\,800 किलोमीटर) कि जवाब उसी चुनाव पर निर्भर करता है।
\end{solution}

\begin{solution}{exo:g4:world-physical-map:9}
3 करोड़ यानी 300 लाख। उसका एक-चौथाई 75 लाख वर्ग किलोमीटर है (300 को 4 से
भाग); एक-तिहाई 100 लाख, यानी 1 करोड़ है (300 को 3 से भाग)। पचासी लाख वाला
सहारा अफ़्रीका के एक-चौथाई से ज़्यादा है, पर एक-तिहाई से
कम।
\end{solution}

\begin{solution}{exo:g4:world-physical-map:10}
$2500 - 1000 = 1500$: झोंपड़ी \omterm{def:g1:landscapes-around-us:water}{झील} से 1\,500 मीटर ऊँची है। 2\,000 और 4\,000 मीटर के
बीच की झोंपड़ी भूरी है; 1\,000 मीटर पर \omterm{def:g1:landscapes-around-us:water}{झील} हल्के पीले (500 से 1\,000) और हल्के
भूरे (1\,000 से 2\,000) की सीमा पर है: दोनों में से कोई भी जवाब
माना जाएगा।
\end{solution}

\begin{solution}{exo:g4:world-physical-map:11}
10 किलोमीटर 10\,000 मीटर हैं। $10000 - 8849 = 1151$: हवाई जहाज़ माउंट एवरेस्ट से
1\,151 मीटर ऊँचा है। चैलेंजर गर्त के ऊपर वह \omterm{def:g2:water-on-earth:ocean}{समुद्र} से 10\,000 मीटर ऊपर है, और
गर्त का तल \omterm{def:g2:water-on-earth:ocean}{समुद्र} से 11\,000 मीटर नीचे है: गर्त के तल से
$10000 + 11000 = 21000$ मीटर
ऊपर।
\end{solution}
